% Pursuing Stacks (cotes 134-2 to 134-8): archive cote 134-4, site page 60 (PDF page 61), typescript page 148.
% typescript. Draft compared: ps3.tex:3053-3116.
% AI-transcribed by Claude Opus 5.5 (claude-opus-5-5, adaptive thinking, effort high), Message Batches API, batch ps2 (msgbatch_01NVwr1Yuo8qJf5rykeGBXPK), 2026-11-12; not yet checked by eye.
\scanpage{134-4/60}
\mspage{148}

What we \uline{can} do however at present,  in terms of a given functor $\hand{\pi}_0$ ,
is to introduce the corresponding notion of 0-connectedness (understood:
with  respect to $\hand{\pi}_0$)  namely objects such that $\hand{\pi}_0$(X) is a one-point set,
the notion of compatibility of a homotopy structure h\xout{mm} on M with $\hand{\pi}_0$ ,
and the $\hand{\pi}_0$-canonical (or simply, ``canonical'') homotopy structure on M,
which now should be denoted by  $h_{\hand{\pi}_0}$  rather than $h_M$  (unless we write $h_{\underline{M}}$
where $\underline{M}$  denotes the pair $(M,\hand{\pi}_0)$), which is the widest weak homotopy in-%
terval structure on M which is $\hand{\pi}_0$-admissible,  and can be described \inserted{(assuming}
$\hand{\pi}_0$ commutes with \inserted{finite} products) \unsure{\over{in}{by}|in} terms of all 0-connected intervals as a
generating family of weak homotopy intervals.  The generalities of par.54
about relationship of $h_{\underline{M}}$ with $h_W$  (where $W\subset F\hand{\ell}(M)$)  should carry over
verbatim,  as well as those of the next\inserted{, \unsure{one},}  provided everywhere  0-con-%
nectedness is understood relative to the given functor $\hand{\pi}_0$ ,  and ``total
0-connectedness''  is interpreted as just meaning that $\hand{\pi}_0$ commutes to fi-%
nite products.  Thus our contact with ``geometry'' via true honest connec\xout{\illegible{1}}%
ted components of objects was of short duration,  and back we are to pure
algebra with just a functor given which we call $\hand{\pi}_0$ God knows why  -  the
culprit for this change of perspective being p\unsure{\over{p}{o}}or modelizer (Spaces),
who was suppose\unsure{\over{\illegible{1}}{d}} to represent the tie with so-called ``topology'' ...

\margin{60.}  I almost forgot I still have to check ``handyness'' of the notions
developped yesterday,  on the example of test categories or rather,  the
corresponding elementary modelizers A\^{} .  As usual,  I can't resist being
a little more general,  so let's start with an arbitrary topos $\hand{\mathcal{A}}$ first.
It always satisfies conditions a) and b) of page 123.  Condition c),  na-%
mely that every\xout{\illegible{1}} object of $\hand{\mathcal{A}}$ could be decomposed into a sum of 0-connec%
-ted ones,  is equivalent with saying that $\hand{\mathcal{A}}$ admits a generating subcate-%
gory  A \unsure{\hand{$\circ$}}made up with 0-connected objects.  In this case, $\hand{\mathcal{A}}$ is called
\uline{locally} 0-\uline{connected} or simply,  \uline{locally connected}  -  which generalizes
the notion known under this name from topological spaces to topoi. \inserted{On the other hand,} \unsure{\over{The}{the}}
condition d) is expressed by saying that the topos considered is \inserted{0-}\uline{connected}
- equally a generalization of the corresponding notion for spaces.  Condi-%
tion e) \deleted{\unsure{however}},  about the product of two 0-connected objects
being 0-connected,  is a highly unusual one in ordinary topology.  For a
topological space,  it means that the space is irreducible (hence reduced
to a point if the space is Hausdorff).  In accordance with the terminology
introduced yesterday,  we'll say that $\hand{\mathcal{A}}$ is \uline{totally} 0-\uline{connected} if it is
locally connected,  and if the product of two 0-connected objects is again
0-connected.  The standart arguments show that for this,  it is enough

\margin{60}
